\section{模型评价与改进方向}

\subsection{模型优点}

我们围绕算力-电力协同调度问题，建立了从基础调度到全要素集成的四层递进的模型体系，在建模方法、约束刻画与可解性设计上具有以下优点：

1. 问题一之所以选用极差作为公平性度量，是因为其数学形式简洁，实际含义明确，直接对应最繁忙与最空闲的区域负载差距，相比方差有更好的可读性与决策效率；仅基于24点预测即实现25.2\%的极差改善。

2. 问题二利用碳排放与成本的同向性（高碳区通常高电价），使碳约束与成本优化方向部分一致，规避多目标冲突的无解困境；依赖滚动时域MILP约束矩阵稀疏特质，我们借助贪心初始解可在数秒内完成单段优化，具备实际部署条件。

3. 问题三通过线性耦合构造充放电与购电时间的约束，把储能充放电省下的钱和优化购电时间省下的钱分开计算，结果分别得到99.36 M与75.15 M元，LP形式保证全局最优。

4. 问题四需要处理六个区域、多个时段，以及储能状态等大量变量。我们采用基荷预填策略来应对高维MILP挑战，预先固定必须运行的任务，将整数规模压缩约$10^5$倍。至此，原本不可解的实例得以在数秒内求解。该预填方案的EDF次优性经Clean-test量化，小于0.88\%。对$\theta=0$全任务MILP这一极端情形，我们另以大邻域搜索（LNS）在50轮/区内收敛作为验证通道，较EDF基线的次优间隙全程约2.3\%。该策略兼顾效率与解质量，是本文大规模MILP求解的关键设计。

\subsection{模型局限}

在肯定上述优点的同时，我们的模型也存在以下局限：

1. 问题一只用未来24点做预测，无法感知更长周期的资源负载趋势。而且跨预测时段边界时衔接生硬，方差目标下最优解不唯一，等效分配的实际负载差异可达3\%--5\%。需通过极差筛选确定可执行方案。

2. 问题二贪心分配对遍历顺序敏感，这一问题尚未彻底解决。不同排序准则可能使分配偏离全局最优。同时贪心阶段以均价近似峰谷平电价，峰谷价差大时引入的分配偏差在后续MILP中仅能部分修正，初始解质量限制了更优解的搜索。

3. 问题三的受限消纳假设实际上主导结果形态，储能与购电的协同收益，依赖消纳上限等同于实际新能源出力的假设。消纳上限放松时净购电标准差反升。这说明储能价值部分源于该假设，结论不应脱离该假设向外推广。

4. 问题四的基荷预填本质是施加额外约束。基荷任务被强制固定后失去与其他任务的交互优化机会。主解在F区存在1小时轻微容量越限及约0.88\%的最优性差距，均与基荷的结构刚性有关。同时，EDF不感知实时电价波动，电价剧烈变化时基荷预填可能做出错误判断。

\subsection{改进方向}

针对上述局限，我们提出以下具体改进方向：

1. 问题一可尝试引入滚动时域控制消除时域边界效应，并以极差作二次惩罚项锁定唯一的可执行解。

2. 问题二可改用拉格朗日松弛分解替换贪心，消除遍历顺序敏感性，同时以分时电价替代均价近似。

3. 问题三考虑将消纳上限扩展为灵敏度变量，画出碳排放变动对成本的响应曲线。然后引入随机规划增强对预测误差的抵抗能力。

4. 问题四可以将基荷预填由硬固定改为软惩罚，不再强制开工时间，转向给偏离预填方案的安排增加代价。同时为大邻域搜索引入自适应邻域半径，优先修复容量紧张的区域。